% titel: Blatt 3, A2: Reaktionsgleichungen ausgleichen
% Teil II: der durchgerechnete Weg.

\subsection*{Bestimmung}
Alle fünf Gleichungen werden so ausgeglichen, dass pro Element (und bei d) auch
die Ladung) links und rechts gleich viele Teilchen stehen.

\subsection*{Rechnung}
\begin{itemize}
  \item[a)] $\mathrm{TiO_2} + \mathrm{Cl_2} + \mathrm{C}
        \rightarrow \mathrm{TiCl_4} + \mathrm{CO}$: \quad Ti als erstes, dann
        O und Cl koppeln über den Faktor $2$ vor $\mathrm{C}$ und $\mathrm{Cl_2}$:
        $2\,\mathrm{O}$ links $\Rightarrow$ $2\,\mathrm{CO}$; $4\,\mathrm{Cl}$
        links $\Rightarrow$ $2\,\mathrm{Cl_2}$.
        \[
          \mathrm{TiO_2} + 2\,\mathrm{Cl_2} + 2\,\mathrm{C}
          \ \rightarrow\ \mathrm{TiCl_4} + 2\,\mathrm{CO}.
        \]
        Prüfen: Ti $1/1$, O $2/2$, Cl $4/4$, C $2/2$.
  \item[b)] $4\,\mathrm{NH_3} + 5\,\mathrm{O_2}
        \rightarrow 4\,\mathrm{NO} + 6\,\mathrm{H_2O}$:
        Saure Stickstoffstickstoffgleichung aus $n\,\mathrm{NH_3}$ rechts
        $\rightarrow n\,\mathrm{NO}$; Wasserstoff aus $3n$ links $\rightarrow$
        $3n/2\,\mathrm{H_2O}$; das kleinste ganzzahlige $n=4$; dann Sauerstoff:
        $2\cdot 4 = 4 + 6 = 10$, also $5\,\mathrm{O_2}$.
        \[
          4\,\mathrm{NH_3} + 5\,\mathrm{O_2}
          \ \rightarrow\ 4\,\mathrm{NO} + 6\,\mathrm{H_2O}.
        \]
        Prüfen: N $4/4$, H $12/12$, O $10/10$.
  \item[c)] $3\,\mathrm{CCl_4} + 2\,\mathrm{SbF_3}
        \rightarrow 3\,\mathrm{CCl_2F_2} + 2\,\mathrm{SbCl_3}$:
        Halogen-Austausch; C und Sb treten je einmal auf, also Faktoren $3$ und
        $2$, damit F und Cl ausgehen:
        \[
          3\,\mathrm{CCl_4} + 2\,\mathrm{SbF_3}
          \ \rightarrow\ 3\,\mathrm{CCl_2F_2} + 2\,\mathrm{SbCl_3}.
        \]
        Prüfen: C $3/3$, Cl $12/(6+6)=12$, Sb $2/2$, F $6/6$.
  \item[d)] $2\,\mathrm{S_2O_3^{2-}} + \mathrm{I_2}
        \rightarrow 2\,\mathrm{I^-} + \mathrm{S_4O_6^{2-}}$:
        Redox) das geschwefelte Ion dimerisiert; $2\,\mathrm{S_2O_3^{2-}}$ gibt
        $\mathrm{S_4O_6^{2-}}$, das Iod wird $\mathrm{I_2} \rightarrow
        2\,\mathrm{I^-}$ (Ladung $2\cdot 2^-$ links $\rightarrow 2\cdot 1^-$
        rechts).
        \[
          2\,\mathrm{S_2O_3^{2-}} + \mathrm{I_2}
          \ \rightarrow\ 2\,\mathrm{I^-} + \mathrm{S_4O_6^{2-}}.
        \]
        Prüfen: S $4/4$, O $6/6$, I $2/2$, Ladung $-4/-4$.
  \item[e)] Setzt man $a$ Formeleinheiten
        $\mathrm{Ca_3(PO_4)_2}$ ein, so sind $3a$ Calcium-Atome zu
        $\mathrm{CaSiO_3}$ zu bilanzieren: $3a\,\mathrm{SiO_2}$ liefern das
        Silicium. Die $4a$ Phosphor-Atome bilden $a\,\mathrm{P_4}$. Der
        Sauerstoff der Phosphate und der Silicate muss in $\mathrm{CO}$
        enden. Mit $a=2$:
        \begin{align*}
          2\,\mathrm{Ca_3(PO_4)_2} &: \mathrm{Ca}_6,\ \mathrm{P}_4,\ \mathrm{O}_{16},\\
          6\,\mathrm{SiO_2} &: \mathrm{Si}_6,\ \mathrm{O}_{12} \quad
          \Longrightarrow \ \mathrm{Summe}\ \mathrm{O}_{28}.
        \end{align*}
        Die $6\,\mathrm{CaSiO_3}$ nehmen $6\cdot 3 = 18\,\mathrm{O}$ auf.
        Verbleiben $28 - 18 = 10\,\mathrm{O}$
        für $\mathrm{CO}$; das fordert $10\,\mathrm{C}$:
        \[
          2\,\mathrm{Ca_3(PO_4)_2} + 6\,\mathrm{SiO_2} + 10\,\mathrm{C}
          \ \rightarrow\ 6\,\mathrm{CaSiO_3} + 10\,\mathrm{CO} + \mathrm{P_4}.
        \]
        Prüfen: Ca $6/6$, P $4/4$, Si $6/6$, O $16+12=28 / 18+10=28$,
        C $10/10$.
\end{itemize}

\begin{ergebnis}
\begin{gather*}
  \mathrm{TiO_2} + 2\,\mathrm{Cl_2} + 2\,\mathrm{C}
    \rightarrow \mathrm{TiCl_4} + 2\,\mathrm{CO}, \\
  4\,\mathrm{NH_3} + 5\,\mathrm{O_2}
    \rightarrow 4\,\mathrm{NO} + 6\,\mathrm{H_2O}, \\
  3\,\mathrm{CCl_4} + 2\,\mathrm{SbF_3}
    \rightarrow 3\,\mathrm{CCl_2F_2} + 2\,\mathrm{SbCl_3}, \\
  2\,\mathrm{S_2O_3^{2-}} + \mathrm{I_2}
    \rightarrow 2\,\mathrm{I^-} + \mathrm{S_4O_6^{2-}}, \\
  2\,\mathrm{Ca_3(PO_4)_2} + 6\,\mathrm{SiO_2} + 10\,\mathrm{C}
    \rightarrow 6\,\mathrm{CaSiO_3} + 10\,\mathrm{CO} + \mathrm{P_4}.
\end{gather*}
\end{ergebnis}

\begin{kontrolle}
Jede Zeile elementweise geprüft (in der Rechnung vermerkt): Atomanzahl und
Ladung stimmen beidseitig überein. Bei a) stimmen $2\,\mathrm{Cl_2}$ mit
$4\,\mathrm{Cl}$ in $\mathrm{TiCl_4}$ überein; bei d) kompensieren die beiden
$\mathrm{S_2O_3^{2-}}$ die $2\,\mathrm{I^-}$-Ladung exakt. Die Koeffizienten
sind die kleinsten ganzzahligen (ggT $=1$).
\end{kontrolle}

\begin{merksatz}
Ausgleichen = Elementbilanz + Ladungsbilanz. Erst Metalle, dann Nichtmetalle
(ohne O, H), zuletzt Sauerstoff und Wasserstoff; Redoxgleichungen über die
Halbreaktionen-Methode.
\end{merksatz}
